\newif\ifglobal

%%%%%%%%%%%%%%%%%%%%%%%
%%% Compile to view %%%
%%%%%%%%%%%%%%%%%%%%%%%

%%%%%%%%%%%%%%%%%%%%%%%%%%%%%%%%%%%%%%%%%%%%%%%%%%%
%%% Readme:                                     %%%
%%% https://www.overleaf.com/read/bwdjvfjksvjy/ %%%
%%%%%%%%%%%%%%%%%%%%%%%%%%%%%%%%%%%%%%%%%%%%%%%%%%%

% >>> M?: marks a module setting number or important notice

% >>> M4: Set {./relative-path} to external files for this module <<<
\def \modulefiles {./27-RNG}
% To add an external file, use:
% (Replace '---' with actual filename)
% '\input{\modulefiles/tables/---.tex}' for tables
% '\includegraphics{\modulefiles/figures/---}' for figures
% '\subfile{\modulefiles/---}' for register files

% >>> M1: This module is in global project? <<<
% 'YES' - uncomment; 'NO' - comment out
\globaltrue

\ifglobal
    \documentclass[main\_\_EN.tex]{subfiles}

\else
    \documentclass[openany, 10pt]{book}

    % >>> M2: In ./00-shared/preamble-trm-module, also find
    %         settings for visibility of todo notes, watermark <<<
    \usepackage{preamble-trm-module}

    % >>> M3: Temporary labels in English,
    %         you can add your own labels there <<<
    \myexternaldocument{./00-shared/temp-labels-en}

\fi %\ifglobal


\setcounter{secnumdepth}{4}


% >>> M5: If this module needs any updates in
% './00-shared/preamble-trm-module.sty'
% (include additional package, etc.), contact your module PM
% or create the file './00-shared/changelog.tex' make the
% necessary updates yourself and document them in this file <<<

% Make text left-aligned, comment out for full justification
\raggedright


\begin{document}


% Sets language into which pre-defined bits of text are translated
\selectlanguage{english}

% Do not delete this block
\ifglobal
% Add the following front matter if
% this module is outside of global project
\else
    \InputIfFileExists{readme}{}{}
    {\let\clearpage\relax \listoftodos}
    {\let\clearpage\relax \listofdonetodos}
    \tableofcontents
    %\listoftables
    %\listoffigures
\fi


%%%%%%%%%%%%%%%%%%%%%%%%%
%%% TRM Module Begins %%%
%%%%%%%%%%%%%%%%%%%%%%%%%

\hypertarget{rng}{}
\chapter{Random Number Generator (RNG)}\label{mod:rng}

\section{Introduction}

The \chipname{} contains a true random number generator, which generates 32-bit random numbers for cryptographical operations and other applications.

\section{Features}

The random number generator in \chipname{} generates true random numbers, derived from a physical process rather than an algorithm. Each number within the specified range has an equal probability of occurrence.

\section{Functional Description}\label{sec:rng-func-descr}

Every 32-bit value that the system reads from the \hyperref[regdesc:LPPERIRNGDATAREG]{LPPERI\_RNG\_DATA\_REG} register of the random number generator is a true random number. These true random numbers are generated based on the \textbf{thermal noise} in the system and the \textbf{asynchronous clock mismatch}.

\begin{itemize}
    \item \textbf{Thermal noise} comes from the high-speed ADC or SAR ADC or both. Whenever the high-speed ADC or SAR ADC is enabled, bit streams will be generated and fed into the random number generator through an XOR logic gate as random seeds.
    \item RC\_FAST\_CLK is an \textbf{asynchronous clock} source and it increases the RNG entropy by introducing circuit metastability.
\end{itemize}


\begin{figure}[ht]
    \centering
    \includegraphics[width=0.8\textwidth]{\modulefiles/figures/29-RNG-noise-source.pdf}
    \caption{Noise Source}
    \label{fig:rng-noise-source}
\end{figure}


When there is noise coming from the high-speed ADC, the random number generator is fed with a 2-bit entropy in one APB clock cycle, which is normally 32 MHz.

When the SAR ADC generates noise, the random number generator obtains 2 bits of entropy per SAR ADC sampling cycle. The SAR ADC sampling cycle is typically 66 kHz, determined by the SAR ADC configuration.

A data sample of 2 GB, which is read from the random number generator at a rate of 5 MHz with only the high-speed ADC being enabled, has been tested using the Dieharder Random Number Testsuite (version 3.31.1). The sample passed all tests.


\section{Programming Procedure}

When using the random number generator, make sure at least either the SAR ADC, high-speed ADC\hyperref[others:rng-note-clock-seed]{$^1$} is enabled and it is recommended to enable  RC\_FAST\_CLK\hyperref[others:rng-note-clock-seed]{$^2$} to increase the entropy which entry random numbers. Otherwise, pseudo-random numbers will be returned.
\begin{itemize}
    \item SAR ADC can be enabled by using the DIG ADC controller. For details, please refer to Chapter \ref{mod:adc} \textit{\nameref{mod:adc}}.\\
    It is recommended to use the following configurations for SAR ADC:
    \begin{itemize}
        \item Set the sampling channel to collect the internal voltage of the chip
        \item Set the controller to use timer-driven continuous sampling mode
    \end{itemize}
    \item High-speed ADC is enabled automatically when the Wireless module is enabled.
    \item RC\_FAST\_CLK is enabled by default and it could not be closed.
\end{itemize}

\begin{tiplisting}\label{others:rng-note-clock-seed}
\vspace{-2.5em}
\begin{enumerate}
    \item Note that, when the Wireless module (Thread or BLE) is enabled, the value read from the high-speed ADC can be saturated in some extreme cases, which lowers the entropy. Thus, it is advisable to also enable the SAR ADC as the noise source for the random number generator for such cases.
    \item Enabling RC\_FAST\_CLK increases the RNG entropy. However, to ensure maximum entropy, it's recommended to always enable an ADC source as well.
\end{enumerate}
\end{tiplisting}

When using the random number generator, read the \hyperref[regdesc:LPPERIRNGDATAREG]{LPPERI\_RNG\_DATA\_REG} register multiple times until sufficient random numbers have been generated.

Please ensure that the rate at which the register is read does not exceed the entropy generation frequency described in section \ref{sec:rng-func-descr} above.


%%%%%%%%%%%%%%%%%%%%%%%%
%%% Register Summary %%%
%%%%%%%%%%%%%%%%%%%%%%%%

\hypertarget{rng-reg-summ}{}
\section{Register Summary}
\label{sec:rng-reg-summ}

The abbreviations given in Column \textbf{Access} are explained in Section \hyperref[glossary-access-types]{Access Types for Registers}.

\begin{table}[H]
    \centering
    \setlength{\extrarowheight}{0.3em}
        \begin{tabular}{ | l | l | l | l | }
    \hline
    \rowcolor{lightgray}
    \textbf{Name} & \textbf{Description} & \textbf{Address} & \textbf{Access} \\ \hline
    \hyperref[regdesc:LPPERIRNGDATAREG]{LPPERI\_RNG\_DATA\_REG} & Random number data & 0x600B\_280C & RO \\ \hline
    \end{tabular}
\end{table}
%%%%%%%%%%%%%%%%%
%%% Registers %%%
%%%%%%%%%%%%%%%%%

\section{Register}

\begin{register}{H}{LPPERI\_RNG\_DATA\_REG}{0x600B\_280C}\label{regdesc:LPPERIRNGDATAREG}
\regfield{LPPERI\_RNG\_DATA}{32}{0}{{0x00000000}}%
\reglabel{Reset}\regnewline%
\begin{regdesc}\begin{reglist}
\item [LPPERI\_RNG\_DATA] Random number source. (RO)
\end{reglist}\end{regdesc}
\end{register}
\end{document}
